\documentclass[10pt, aspectratio=169]{beamer}
\input{../../_en_preambulo_beamer}
% Comandos y macros estadísticas compartidas para las presentaciones Beamer en inglés (Metropolis)

% Conjuntos numéricos básicos
\newcommand{\R}{\mathbb{R}}
\newcommand{\N}{\mathbb{N}}
\newcommand{\Q}{\mathbb{Q}}
\newcommand{\Z}{\mathbb{Z}}
\newcommand{\set}[1]{\left\{ #1 \right\}}
\newcommand{\abs}[1]{\left|#1\right|}
\newcommand{\norm}[1]{\left\| #1 \right\|}

% Comandos estadísticos y probabilísticos del curso
\newcommand{\Var}{\operatorname{Var}}
\newcommand{\cov}{\operatorname{Cov}}
\newcommand{\corr}{\rho}
\newcommand{\comb}[2]{\begin{pmatrix} #1 \\ #2 \end{pmatrix}}
\newcommand{\s}{\sigma}
\newcommand{\particion}{\mathfrak{P}}
\newcommand{\card}[1]{\#\left( #1 \right)}
\newcommand{\seq}[3]{\set{#1}_{#2}^{#3}}
\newcommand{\E}{\mathbb{E}}
\newcommand{\Prob}{\mathbb{P}}

% Entornos pedagógicos y cajas en inglés académico natural (soportando tanto nombres en inglés como en español)
\newenvironment{discussion}{%
  \begin{alertblock}{Discussion Question}%
}{%
  \end{alertblock}%
}
\newenvironment{discusion}{%
  \begin{alertblock}{Discussion Question}%
}{%
  \end{alertblock}%
}

\newenvironment{reflection}{%
  \begin{alertblock}{Classroom Reflection}%
}{%
  \end{alertblock}%
}
\newenvironment{reflexion}{%
  \begin{alertblock}{Classroom Reflection}%
}{%
  \end{alertblock}%
}

\newenvironment{paradox}{%
  \begin{alertblock}{Exploratory Paradox}%
}{%
  \end{alertblock}%
}
\newenvironment{paradoja}{%
  \begin{alertblock}{Exploratory Paradox}%
}{%
  \end{alertblock}%
}

\newenvironment{motivation}{%
  \begin{exampleblock}{Motivation: Data Science in Practice}%
}{%
  \end{exampleblock}%
}
\newenvironment{motivacion}{%
  \begin{exampleblock}{Motivation: Data Science in Practice}%
}{%
  \end{exampleblock}%
}

\newenvironment{quickchallenge}{%
  \begin{block}{Quick Team Challenge}%
}{%
  \end{block}%
}
\newenvironment{retoclase}{%
  \begin{block}{Quick Team Challenge}%
}{%
  \end{block}%
}


\title[Multiple Regression]{Section 09.02: Multiple Linear Regression}
\subtitle{The Normal Equation in Matrix Notation and Ridge/Lasso Regularization}
\author[J. Castillo Colmenares]{Juliho Castillo Colmenares}
\institute{Tecnológico de Monterrey}
\date{\vspace{-1.5cm}}

\begin{document}

% Slide 1: Title
\begin{frame}[plain]
  \titlepage
\end{frame}

% Slide 2: Roadmap
\begin{frame}{Roadmap --- Chapter 09: Linear and Multiple Regressions}
  \footnotesize
  Where are we in the modular development of this chapter? \pause
  \begin{itemize}\itemsep=0.08em
    \item \textbf{09.01} Simple Linear Regression (OLS) and $R^2$ \pause
    \item \textbf{09.02} Multiple Linear Regression, the Normal Equation, and Ridge/Lasso \emph{(Today)} \pause
    \item \textbf{09.03} Residual Diagnostics, Multicollinearity (VIF), and Assumptions \pause
    \item \textbf{09.04} Model Validation, $k$-fold Cross-Validation, and \texttt{scikit-learn}
  \end{itemize}
  \vspace{-0.05cm}
  \begin{block}{Session Objective}
    Generalize simple regression to $p$ simultaneous predictors via matrix notation: derive the Normal Equation $\hat{\boldsymbol\beta}=(\mathbf{X}^T\mathbf{X})^{-1}\mathbf{X}^T\mathbf{Y}$, and present Ridge and Lasso regularization as a remedy when the design matrix is close to singular due to multicollinearity.
  \end{block}
\end{frame}

% Slide 3: Motivation
\begin{frame}{From One Predictor to $p$ Simultaneous Predictors}
  \footnotesize
  \begin{motivacion}
    In practice a single variable rarely explains a phenomenon: a vehicle's fuel efficiency depends on its weight \emph{and} its engine size; salary depends on experience \emph{and} education level. Writing and solving the normal equations with scalar summations for each $\beta_j$ becomes unmanageable as $p$ grows.
  \end{motivacion}

  \vspace{0.15cm}
  \pause
  \begin{itemize}\itemsep=0.1cm
    \item Matrix notation solves the problem compactly and exactly for any $p$. \pause
    \item The model, the objective function, and the solution have the \textbf{same structure} as in simple regression, only expressed with vectors and matrices.
  \end{itemize}
\end{frame}

% Slide 4: Matrix model
\begin{frame}{The Multiple Linear Model in Matrix Form}
  \footnotesize
  \begin{block}{Population Model}
    $$\mathbf{Y}=\mathbf{X}\boldsymbol\beta+\boldsymbol\varepsilon, \qquad E[\boldsymbol\varepsilon]=\mathbf{0},\ \text{Var}(\boldsymbol\varepsilon)=\sigma^2\mathbf{I}_n$$
  \end{block}
  \pause

  \vspace{0.1cm}
  \begin{alertblock}{Dimensions}
    $\mathbf{Y}$ is $n\times1$; $\mathbf{X}$ is $n\times(p+1)$ (first column of ones); $\boldsymbol\beta$ is $(p+1)\times1$. The objective function is the quadratic form $S(\boldsymbol\beta)=(\mathbf{Y}-\mathbf{X}\boldsymbol\beta)^T(\mathbf{Y}-\mathbf{X}\boldsymbol\beta)$.
  \end{alertblock}
\end{frame}

% Slide 5: Normal equation
\begin{frame}{Gauss's Normal Equation}
  \footnotesize
  \begin{block}{Closed-Form Solution}
    $$\mathbf{X}^T\mathbf{X}\hat{\boldsymbol\beta}=\mathbf{X}^T\mathbf{Y} \implies \hat{\boldsymbol\beta}=(\mathbf{X}^T\mathbf{X})^{-1}\mathbf{X}^T\mathbf{Y}$$
  \end{block}
  \pause

  \vspace{0.1cm}
  \begin{alertblock}{Hat Matrix}
    $\mathbf{H}=\mathbf{X}(\mathbf{X}^T\mathbf{X})^{-1}\mathbf{X}^T$ projects $\mathbf{Y}$ onto $\hat{\mathbf{Y}}=\mathbf{H}\mathbf{Y}$. It is symmetric ($\mathbf{H}^T=\mathbf{H}$) and idempotent ($\mathbf{H}^2=\mathbf{H}$): an exact orthogonal projection onto the column space of $\mathbf{X}$.
  \end{alertblock}
\end{frame}

% Slide 6: Regularization motivation
\begin{frame}{When $\mathbf{X}^T\mathbf{X}$ Is Nearly Singular}
  \footnotesize
  \begin{block}{The Multicollinearity Problem}
    If two or more predictors are highly correlated, $\mathbf{X}^T\mathbf{X}$ approaches singularity: its inversion becomes numerically unstable and the variances of $\hat{\boldsymbol\beta}$ explode (formally diagnosed via \texttt{VIF} in Section 09.03).
  \end{block}
  \pause

  \vspace{0.1cm}
  \begin{alertblock}{The Idea Behind Regularization}
    Penalizing the magnitude of the coefficients in the objective function stabilizes the solution in exchange for introducing controlled bias --- the \textbf{bias-variance tradeoff}.
  \end{alertblock}
\end{frame}

% Slide 7: Ridge and Lasso
\begin{frame}{Ridge ($L_2$) and Lasso ($L_1$) Regression}
  \footnotesize
  \begin{block}{Ridge}
    $$S_{\text{Ridge}}=\|\mathbf{Y}-\mathbf{X}\boldsymbol\beta\|^2+\lambda\|\boldsymbol\beta\|_2^2 \implies \hat{\boldsymbol\beta}_{\text{Ridge}}=(\mathbf{X}^T\mathbf{X}+\lambda\mathbf{I})^{-1}\mathbf{X}^T\mathbf{Y}$$
  \end{block}
  \pause

  \vspace{0.1cm}
  \begin{alertblock}{Lasso}
    $$S_{\text{Lasso}}=\|\mathbf{Y}-\mathbf{X}\boldsymbol\beta\|^2+\lambda\|\boldsymbol\beta\|_1 \quad \text{(no closed form; induces \emph{sparsity})}$$
  \end{alertblock}
\end{frame}

% Slide 8: Python Lab (Block 1)
\begin{frame}[fragile]{Python Lab: Normal Equation and Hat Matrix}
  \scriptsize
  Matrix solution $\hat{\boldsymbol\beta}$ verified against \texttt{np.linalg.lstsq}, and numerical verification of the symmetry/idempotence of $\mathbf{H}$ (\texttt{09.02\_multiple\_linear\_regression.py}):
  {\lstset{basicstyle=\fontsize{3.4pt}{4.1pt}\selectfont\ttfamily,aboveskip=0.05em,belowskip=0.05em}
  \lstinputlisting[firstline=18, lastline=40]{../../code/09_regresiones/09.02_multiple_linear_regression.py}}
\end{frame}

% Slide 9: Python Lab (Block 2)
\begin{frame}[fragile]{Python Lab: Ridge Regression}
  \scriptsize
  Closed-form Ridge solution for increasing $\lambda$, observing coefficient shrinkage:
  {\lstset{basicstyle=\fontsize{3.8pt}{4.5pt}\selectfont\ttfamily,aboveskip=0.05em,belowskip=0.05em}
  \lstinputlisting[firstline=43, lastline=51]{../../code/09_regresiones/09.02_multiple_linear_regression.py}}
\end{frame}

% Slide 10: Python Lab (Block 3)
\begin{frame}[fragile]{Python Lab: Lasso via Coordinate Descent}
  \scriptsize
  Manual implementation of the soft-thresholding operator for Lasso:
  {\lstset{basicstyle=\fontsize{3.2pt}{3.9pt}\selectfont\ttfamily,aboveskip=0.05em,belowskip=0.05em}
  \lstinputlisting[firstline=54, lastline=76]{../../code/09_regresiones/09.02_multiple_linear_regression.py}}
\end{frame}

% Slide 11: Terminal Output
\begin{frame}[fragile]{Python Lab: Terminal Output}
  \scriptsize
  Standard output produced when running the lab script:
  \vspace{0.02cm}
  \begin{block}{Standard Console Output}
    \fontsize{3.6pt}{4.3pt}\selectfont
    \begin{verbatim}
=== Block 1: Normal Equation and Hat Matrix ===
beta_hat = [3.1001, 2.1533, -0.5180] (matches np.linalg.lstsq)
Symmetry error=3.26e-16, Idempotence error=4.09e-16

=== Block 2: Ridge Shrinkage ===
lambda=0: [3.10, 2.15, -0.52]  lambda=100: [0.09, 1.41, 0.10]

=== Block 3: Lasso Sparsity ===
lambda=0:    zero coefficients=0/2
lambda=200:  zero coefficients=1/2 (x2 dropped)
lambda=1000: zero coefficients=2/2 (only intercept survives)
    \end{verbatim}
  \end{block}
\end{frame}

% Slide 12A: Exercise Level 1 (Statement)
\begin{frame}{In-Class Exercise --- Level 1: Fundamental (1/2)}
  \footnotesize
  \begin{block}{Problem (Statement, Problem 9.10.1)}
    With $n=30$ vehicles: $\hat Y=25.0-3.0X_1-1.5X_2$, $R^2=0.780$. Interpret $\hat\beta_1,\hat\beta_2$ \emph{ceteris paribus}, predict fuel efficiency for $X_1=1.5$ tons, $X_2=2.0$ L, and compute SSR/SSE given $S_Y^2=12.0$.
  \end{block}

  \vspace{0.15cm}
  \pause
  \begin{alertblock}{Setup}
    $\text{SST}=(n-1)S_Y^2$; $\text{SSR}=R^2\cdot\text{SST}$; $\text{SSE}=\text{SST}-\text{SSR}$.
  \end{alertblock}
\end{frame}

% Slide 12B: Exercise Level 1 (Resolution)
\begin{frame}{In-Class Exercise --- Level 1: Fundamental (2/2)}
  \scriptsize
  We substitute the given values into the fitted equation: \pause

  \vspace{0.05cm}
  \begin{block}{Calculation}
    $$\hat Y=25.0-3.0(1.5)-1.5(2.0)=25.0-4.5-3.0=17.5 \text{ km/l}$$
    $$\text{SST}=(29)(12.0)=348.0, \quad \text{SSR}=0.780(348.0)\approx271.44, \quad \text{SSE}\approx76.56$$
  \end{block}

  \vspace{0.05cm}
  \pause
  \begin{alertblock}{Interpretation}
    Holding $X_2$ fixed, each additional ton of weight reduces efficiency by $3.0$ km/l; holding $X_1$ fixed, each additional liter of engine size reduces it by $1.5$ km/l. The model explains $78\%$ of the total variability in efficiency.
  \end{alertblock}
\end{frame}

% Slide 13A: Exercise Level 2 (Statement)
\begin{frame}{In-Class Exercise --- Level 2: Operational (1/2)}
  \footnotesize
  \begin{block}{Problem --- Partial $F$-Test (Statement, Problem 9.10.4)}
    Model 1 (restricted, $p_1=2$): $R_1^2=0.750$. Model 2 (full, $p_2=3$): $R_2^2=0.780$, $n=40$. Run the Partial $F$-Test ($H_0:\beta_3=0$) at $\alpha=0.05$ ($F_{0.05,1,36}=4.113$).
  \end{block}

  \vspace{0.15cm}
  \pause
  \begin{alertblock}{Strategy}
    \scriptsize
    $F_{\text{partial}}=\dfrac{(R_2^2-R_1^2)/(p_2-p_1)}{(1-R_2^2)/(n-p_2-1)}$.
  \end{alertblock}
\end{frame}

% Slide 13B: Exercise Level 2 (Resolution)
\begin{frame}{In-Class Exercise --- Level 2: Operational (2/2)}
  \scriptsize
  We substitute directly into the incremental statistic formula: \pause

  \vspace{0.05cm}
  \begin{block}{Calculation}
    $$F_{\text{partial}}=\frac{(0.780-0.750)/(3-2)}{(1-0.780)/(40-3-1)}=\frac{0.030/1}{0.220/36}=\frac{0.030}{0.006111}\approx4.909$$
  \end{block}

  \vspace{0.05cm}
  \pause
  \begin{alertblock}{Interpretation}
    Since $4.909>4.113$, \textbf{we reject $H_0$}: adding the third predictor ($X_3$, number of sales staff) significantly improves the fit; Model 2 is preferable to Model 1.
  \end{alertblock}
\end{frame}

% Slide 14A: Exercise Level 3 (Statement)
\begin{frame}{In-Class Exercise --- Level 3: Analytical (1/2)}
  \footnotesize
  \begin{block}{Problem --- Exact Matrix Solution (Statement, Problem 9.10.7)}
    With $n=5$, $\mathbf{X}^T\mathbf{X}=\begin{pmatrix}5&10&20\\10&30&40\\20&40&90\end{pmatrix}$, $\mathbf{X}^T\mathbf{Y}=\begin{pmatrix}50\\120\\230\end{pmatrix}$, and $(\mathbf{X}^T\mathbf{X})^{-1}=\begin{pmatrix}2.2&-0.2&-0.4\\-0.2&0.1&0.0\\-0.4&0.0&0.1\end{pmatrix}$. Compute $\hat{\boldsymbol\beta}$.
  \end{block}

  \vspace{0.1cm}
  \pause
  \begin{alertblock}{Strategy}
    \scriptsize
    $\hat{\boldsymbol\beta}=(\mathbf{X}^T\mathbf{X})^{-1}\mathbf{X}^T\mathbf{Y}$: multiply the inverse matrix by the vector $\mathbf{X}^T\mathbf{Y}$.
  \end{alertblock}
\end{frame}

% Slide 14B: Exercise Level 3 (Resolution)
\begin{frame}{In-Class Exercise --- Level 3: Analytical (2/2)}
  \scriptsize
  We multiply row by column: \pause

  \vspace{0.05cm}
  \begin{block}{Calculation}
    \scriptsize
    $$\hat\beta_0=2.2(50)-0.2(120)-0.4(230)=110-24-92=-6.0$$
    $$\hat\beta_1=-0.2(50)+0.1(120)+0.0(230)=-10+12+0=2.0$$
    $$\hat\beta_2=-0.4(50)+0.0(120)+0.1(230)=-20+0+23=3.0$$
  \end{block}

  \vspace{0.05cm}
  \pause
  \begin{alertblock}{Interpretation}
    $\hat Y=-6.0+2.0X_1+3.0X_2$. This exact miniature calculation ($n=5$) illustrates the same matrix procedure that solves models with hundreds of predictors.
  \end{alertblock}
\end{frame}

% Slide 15A: Exercise Level 4 (Statement)
\begin{frame}{In-Class Exercise --- Level 4: Challenging (1/2)}
  \footnotesize
  \begin{block}{Problem --- Matrix Derivation of OLS (Statement, Problem 9.10.9)}
    Starting from $S(\boldsymbol\beta)=(\mathbf{Y}-\mathbf{X}\boldsymbol\beta)^T(\mathbf{Y}-\mathbf{X}\boldsymbol\beta)$, derive the normal equations $\mathbf{X}^T\mathbf{X}\hat{\boldsymbol\beta}=\mathbf{X}^T\mathbf{Y}$ and prove that $E[\hat{\boldsymbol\beta}]=\boldsymbol\beta$.
  \end{block}

  \vspace{0.1cm}
  \pause
  \begin{alertblock}{Strategy}
    \scriptsize
    Expand $S(\boldsymbol\beta)$, take the vector gradient $\nabla_{\boldsymbol\beta}S=\mathbf{0}$, and substitute the true model $\mathbf{Y}=\mathbf{X}\boldsymbol\beta+\boldsymbol\varepsilon$ into $\hat{\boldsymbol\beta}$.
  \end{alertblock}
\end{frame}

% Slide 15B: Exercise Level 4 (Resolution)
\begin{frame}{In-Class Exercise --- Level 4: Challenging (2/2)}
  \scriptsize
  We expand the quadratic form and differentiate: \pause

  \vspace{0.05cm}
  \begin{block}{Calculation}
    \scriptsize
    $$S(\boldsymbol\beta)=\mathbf{Y}^T\mathbf{Y}-2\boldsymbol\beta^T\mathbf{X}^T\mathbf{Y}+\boldsymbol\beta^T\mathbf{X}^T\mathbf{X}\boldsymbol\beta \implies \nabla_{\boldsymbol\beta}S=-2\mathbf{X}^T\mathbf{Y}+2\mathbf{X}^T\mathbf{X}\boldsymbol\beta=\mathbf{0}$$
    $$\hat{\boldsymbol\beta}=(\mathbf{X}^T\mathbf{X})^{-1}\mathbf{X}^T\mathbf{Y}=(\mathbf{X}^T\mathbf{X})^{-1}\mathbf{X}^T(\mathbf{X}\boldsymbol\beta+\boldsymbol\varepsilon)=\boldsymbol\beta+(\mathbf{X}^T\mathbf{X})^{-1}\mathbf{X}^T\boldsymbol\varepsilon$$
  \end{block}

  \vspace{0.05cm}
  \pause
  \begin{alertblock}{Interpretation}
    $E[\hat{\boldsymbol\beta}]=\boldsymbol\beta+(\mathbf{X}^T\mathbf{X})^{-1}\mathbf{X}^T E[\boldsymbol\varepsilon]=\boldsymbol\beta+\mathbf{0}=\boldsymbol\beta$: the matrix OLS estimator is unbiased for any number of predictors $p$, exactly generalizing the scalar result from Section 09.01.
  \end{alertblock}
\end{frame}

% Slide 16: Closing
\begin{frame}{Section Closing: Multiple Linear Regression}
  \begin{reflexion}
    We have generalized simple regression to $p$ predictors via matrix notation: the Normal Equation $\hat{\boldsymbol\beta}=(\mathbf{X}^T\mathbf{X})^{-1}\mathbf{X}^T\mathbf{Y}$ solves the entire model in a single step, and Ridge/Lasso stabilize that solution when multicollinearity threatens to make the matrix nearly singular.
  \end{reflexion}

  \vspace{0.2cm}
  \pause
  \begin{block}{Key Takeaways}
    \begin{itemize}\itemsep=0.08cm
      \item \textbf{Normal Equation:} $\mathbf{X}^T\mathbf{X}\hat{\boldsymbol\beta}=\mathbf{X}^T\mathbf{Y}$, an exact generalization of Section 09.01's scalar formula.
      \item \textbf{Ridge:} $L_2$ penalty, closed-form solution, shrinks coefficients without eliminating them.
      \item \textbf{Lasso:} $L_1$ penalty, no closed form, induces \emph{sparsity} (automatic variable selection).
    \end{itemize}
  \end{block}
\end{frame}

% Slide 17: Modular Perspective
\begin{frame}{Modular Perspective --- The Next Challenge}
  \scriptsize
  \begin{retoclase}
    With the multiple model fitted, Section 09.03 asks whether it is \textbf{trustworthy}: we will formally audit the classical assumptions (residual normality and homoscedasticity, independence) and quantify multicollinearity via the Variance Inflation Factor (VIF) --- the metric that determines exactly when Ridge or Lasso are needed.
  \end{retoclase}

  \vspace{0.08cm}
  \pause
  \begin{alertblock}{Recommended Preparation}
    \begin{itemize}\itemsep=0.06cm
      \item Review the geometric interpretation of the residuals as the vector $(\mathbf{I}_n-\mathbf{H})\mathbf{Y}$.
      \item Reflect on why a high VIF does not invalidate the model, but does drastically inflate the standard errors of the coefficients.
    \end{itemize}
  \end{alertblock}
\end{frame}

\end{document}
